\chapter{Apr.~13 --- Quantum Groups, Part 2}

\section{Quantum Groups, Continued}

\begin{remark}
  Recall the algebra $U_q(\g)$ defined last time.
  It is a Hopf algebra with coproduct
  \[
    \Delta(q^h)
    = q^h \otimes q^h, \quad
    \Delta(e_i)
    = e_i \otimes q^{d_i h_i} + 1 \otimes e_i,
    \quad \Delta(f_i)
    = f_i \otimes 1 + q^{-d_i h_i} \otimes f_i.
  \]
  with counit $\varepsilon(e_i) = \varepsilon(f_i) = 0$ and $\varepsilon(q^h) = 1$.
  The antipode is given by
  \[
    \gamma(e_i)
    = -e_i q^{-d_i h_i}, \quad
    \gamma(f_i) = -q^{d_i h_i} f_i, \quad
    \gamma(q^h) = q^{-h}.
  \]
  There are two Hopf subalgebras
  $U_q(\g^+)$ and $U_q(\g^-)$ generated by
  $\{q^h, e_i\}$ and $\{q^h, f_i\}$, respectively.
\end{remark}

\section{Representations with Weight Decomposition}

\begin{remark}
  Suppose that we have a weight decomposition
  $V = \bigoplus_\mu V^\mu$, where
  $q^h v = q^{\mu(h)} v$ for $v \in V^\mu$
  and $\mu \in \mathfrak{h}^*$.
  Similarly as before, we have the Verma module
  \[
    M_\lambda^{(q)}
    = \Ind^{U_q(\g)}_{U_q(\g^+)} X_\lambda,
  \]
  where $X_\lambda$ is the $1$-dimensional
  module where all $e_i$ act as $0$.
  We can also define
  $L_\lambda^{(q)} = M_{\lambda}^{(q)} / I_\lambda$,
  where $I_\lambda$ is the maximal proper
  submodule of $M_\lambda^{(q)}$.
\end{remark}

\begin{theorem}
  Let $\g$ be a simple finite-dimensional
  Lie algebra and $q \in \C$ (not a root
  of unity). Then
  \begin{enumerate}
    \item $L_\lambda^{(q)}$ is finite-dimensional
      if and only if $\lambda \in P^+$, and
      in this case,
      $\dim L_\lambda^{(q)} = \dim L_\lambda$
      (the same is true for the dimensions
      of the weight subspaces);
    \item the category $C(\g, q)$ of
      finite-dimensional $U_q(\g)$-modules is
      semisimple.
  \end{enumerate}
\end{theorem}

\begin{remark}
  The contragradient module for
  $M = \bigoplus_\lambda M^\lambda$
  is given by
  the restricted dual
  $M^c = \bigoplus_\lambda (M^\lambda)^*$, where
  we have
  \[
    \langle g v^*, v \rangle
    = \langle v^*, \tau(g) v \rangle,
    \quad v^* \in M^c, v \in M, g \in U_q(\g).
  \]
  Here $\tau(e_i) = q^{d_i h_i} f_i$,
  $\tau(f_i) = e_i q^{-d_i h_i}$, and
  $\tau(q^h) = q^h$. Note that
  $\tau(ab) = \tau(b) \tau(a)$, this is done so
  that $(M_1 \otimes M_2)^c \cong M_1^c \otimes M_2^c$
  as $U_q(\g)$-modules.
  Note that the $U_q(\g)$-morphism
  \[
    M_\lambda^{(q)}
    \longrightarrow (M_\lambda^{(q)})^c
  \]
  is generally not an isomorphism, its
  image is $L_\lambda^{(q)}$.
\end{remark}

\begin{remark}
  Let $\rho \in \mathfrak{h}^*$ such that
  $\llparenthesis \rho, h_i \rrparenthesis = 1$,
  i.e. $\alpha_i(\rho) = d_i$. Then
  \[
    \gamma^2(x) = q^{2 \rho} x q^{-2 \rho}.
  \]
  Let $V^*$ be the dual module of $V$.
  Then a priori, $V^{**}$ has a different
  $U_q(\g_e)$-module structure from $V$.\footnote{Recall that $\g_e$ is the extended Lie algebra with an added derivation $d$, which gives a non-degenerate bilinear form $\llparenthesis \cdot, \cdot \rrparenthesis$.}
\end{remark}

\begin{prop}
  If $V$ is a finite-dimensional module over
  $U_q(\g_e)$ with weight decomposition, then
  \begin{align*}
    \delta_V : V &\longrightarrow V^{**} \\
    v &\longmapsto q^{2 \rho} v
  \end{align*}
  is an isomorphism of
  $U_q(\g_e)$-modules.
\end{prop}

\section{Quasitriangular Structure}

\begin{definition}
  Let $A$ be a Hopf algebra and
  $R = \sum_i a_i \otimes b_i \in A \otimes A$
  invertible
  such that\footnote{Here $R_{i, j}$ denotes $R$ acting in the $i, j$ factors of the tensor product $A \otimes A \otimes A$.}
  \[
    R \Delta(x) = \Delta^{\mathrm{op}}(x) R, \quad
    (\Delta \otimes I) R = R_{1, 3} R_{2, 3},
    \quad (I \otimes \Delta) R = R_{1, 3} R_{1, 2}.
  \]
  Such an $R$ is called
  a \emph{universal $R$-matrix} for $A$,
  and $A$ with a choice of
  $R$ is called \emph{quasitriangular}.
\end{definition}

\begin{remark}
  Note that if $V_1, V_2$ are
  two finite-dimensional representations of
  $A$, then
  $V_1 \otimes V_2 \cong V_2 \otimes V_1$.
  Thus we can define
  $\widecheck{R}_{V_1, V_2} = P_{1, 2} R|_{V_1 \otimes V_2}$,
  where $P_{1, 2}$ is the isomorphism
  $V_1 \otimes V_2 \to V_2 \otimes V_1$
  given by permuting the factors.
\end{remark}

\begin{prop}
  We have the following:
  \[
    (\varepsilon \otimes {\id})(R)
    = ({\id} \otimes \varepsilon)(R) = 1_A,
    \quad
    (\gamma \otimes {\id})(R)
    = (\id \otimes \gamma^{-1})(R) = R^{-1}, \quad
    (\gamma \otimes \gamma)(R) = R.
  \]
  Moreover, $R^{\mathrm{op}} = P(R)$
  (where $P(a \otimes b) = b \otimes a$)
  is a universal $R$-matrix for $A^{\mathrm{op}}$.
\end{prop}

\begin{definition}
  A \emph{braided tensor category}
  is an additive category
  $\mathcal{C}$ with a distinguished
  object $1$, a bifunctor
  $\otimes : \mathcal{C} \times \mathcal{C} \to \mathcal{C}$,
  and isomorphisms
  \begin{gather*}
    \alpha_{V_1, V_2, V_3}
    : (V_1 \otimes V_2) \otimes V_3
    \cong V_1 \otimes (V_2 \otimes V_3), \\
    \lambda_V : 1 \otimes V \cong V, \quad
    \rho_V : V \otimes 1 \cong V, \\
    \sigma_{V_1, V_2} : V_1 \otimes V_2 \cong V_2 \otimes V_1
  \end{gather*}
  satisfying the following properties:
  \begin{enumerate}
    \item \emph{triangle axiom:}
      the following diagram commutes:
      \[
        \begin{tikzcd}
          (X \otimes 1) \otimes Y \ar[rr, "\rho_X \otimes {\id}"]
          \ar[dr, "\alpha_{X, 1, Y}", swap]
          & & X \otimes Y \\
          & X \otimes (1 \otimes Y)
          \ar[ur, "{\id} \otimes \lambda_Y", swap]
        \end{tikzcd}
      \]
    \item \emph{pentagon axiom:}
      all maps
      $((V_1 \otimes V_2) \otimes V_3) \otimes V_4 \to V_1 \otimes (V_2 \otimes (V_3 \otimes V_4))$
      obtained by composing
      $\alpha$ and $\alpha^{-1}$ are equal;
    \item \emph{hexagon axioms:}
      the isomorphisms
      $V_1 \otimes (V_2 \otimes V_3) \to (V_2 \otimes V_3) \otimes V_1$
      given by
      $\sigma_{V_1, V_2 \otimes V_3}$ and
      \[
        \alpha^{-1}
        ({\id_{V_2}} \otimes \delta_{V_1, V_3})
        \alpha (\delta_{V_1, V_2} \otimes {\id_{V_3}}) \alpha^{-1}
      \]
      are equal, and similarly for $\sigma^{-1}$.
  \end{enumerate}
\end{definition}

\begin{theorem}
  Let $A$ be a quasitriangular Hopf algebra.
  The category of representations of
  $A$ is a braided tensor category with
  $1 = \C$, trivial unit morphisms and
  associative morphisms, and
  $\delta_{V_, W} = \widecheck{R}_{V, W}$.
\end{theorem}

\begin{remark}
  In this category, every object
  $V$ has a dual $V^*$, and we have morphisms
  $1 \to V \otimes V^*$ and
  $V^* \otimes V \to 1$. This is known as
  \emph{rigidity}.
\end{remark}

\section{Quantum Yang-Baxter Equation}

\begin{prop}[Quantum Yang-Baxter equation]
  We have $R_{1, 2} R_{1, 3} R_{2, 3} = R_{2, 3} R_{1, 3} R_{1, 2}$.
\end{prop}

\begin{proof}
  We can write
  \[
    R_{1, 2} R_{1, 3} R_{2, 3}
    = R_{1, 2} (\Delta \otimes {\id}) R
    = (\Delta^{\mathrm{op}} \otimes {\id})(R)R_{1, 2}
    = P_{1, 2}(R_{1, 3} R_{2, 3}) R_{1, 2}
    = R_{2, 3} R_{1, 3} R_{1, 2},
  \]
  which is the desired formula.
\end{proof}

\begin{definition}
  Define the \emph{braid group on $N$ strands}
  to be the group $B_N$
  generated by $b_1, \dots, b_{N - 1}$
  with relations
  $b_i b_{i + 1} b_i = b_{i + 1} b_i b_{i + 1}$
  and $b_i b_j = b_j b_i$ for $|i - j| > 1$.
\end{definition}

\begin{theorem}
  Let $A$ be a quasitriangular Hopf algebra
  and $V_1, \dots, V_N$ representations
  of $A$. Define
  \[
    \widecheck{R}_i
    = V_1 \otimes \cdots \otimes V_N
    \longrightarrow
    V_1 \otimes \cdots \otimes V_{i + 1}
    \otimes V_i \otimes \cdots \otimes V_N
  \]
  by $\widecheck{R}_i = P_{i, i + 1} R_{V_i, V_{i + 1}}$.
  Then $\widecheck{R}_i \widecheck{R}_{i + 1} \widecheck{R}_i = \widecheck{R}_{i + 1} \widecheck{R}_i \widecheck{R}_{i + 1}$.
\end{theorem}

\begin{corollary}
  In a braided tensor category, we can define
  isomorphisms
  \[
    \beta^{\pm}_{X, Y, Z}
    : X \otimes (Y \otimes Z)
    \longrightarrow
    Y \otimes (X \otimes Z)
  \]
  by $\beta^+_{X, Y, Z} = \alpha(\sigma_{X, Y} \otimes {\id_Z}) \alpha^{-1}$ and
  $\beta^-_{X, Y, Z} = \alpha(\sigma^{-1}_{X, Y} \otimes {\id_Z}) \alpha^{-1}$.
  Let
  \begin{align*}
    \beta^{\pm}_{1, 2}
    = \beta^{\pm}_{X, Y, Z}
    : X \otimes (Y \otimes (Z \otimes U))
    &\longrightarrow
    Y \otimes (X \otimes (Z \otimes U)), \\
    \beta_{2, 3}^{\pm}
    = {\id} \otimes \beta^{\pm}_{Y, Z, U}
    : X \otimes (Y \otimes (Z \otimes U))
    &\longrightarrow
    X \otimes (Z \otimes (Y \otimes U)).
  \end{align*}
  Then we have
  $\beta^{\pm}_{1, 2} \beta^{\pm}_{2, 3} \beta^{\pm}_{1, 2} = \beta^{\pm}_{2, 3} \beta^{\pm}_{1, 2} \beta^{\pm}_{2, 3}$
  and $\beta^{\pm}_{X, Y, Z} \beta^{\mp}_{Y, X, Z} = {\id}$.
\end{corollary}

\section{Quantum Double}

\begin{remark}
  Let $A$ be a finite-dimensional Hopf algebra.
  Let $\{a_i\}$ a basis for $A$, and
  $\{a^i\}$ the dual basis of $A^*$.
  Define $D(A) = A \otimes (A^*)^{\mathrm{op}}$.
  We want a Hopf algebra structure on $D(A)$
  so that the natural embeddings
  $A \hookrightarrow D(A)$ and
  $(A^*)^{\mathrm{op}} \hookrightarrow D(A)$
  are Hopf algebra morphisms.

  One constructs the \emph{quantum double} (or \emph{Drinfeld double})
  as follows. Let the structure constants of
  \[
    A \otimes A \otimes A \longrightarrow A
  \]
  (note that $m(m \otimes {\id}) = m({\id} \otimes m)$)
  be $m^i_{j, k, \ell}$, i.e.
  $a_j a_k a_\ell = \sum_i m^i_{j, k, \ell} a_i$,
  and the structure constants of
  \[
    A \longrightarrow A \otimes A \otimes A
  \]
  (note that $(\Delta \otimes {\id}) \Delta = ({\id} \otimes \Delta) \Delta$)
  be $\mu^{j, k, \ell}_i$. Then we can define
  multiplication on $D(A)$ by
  relations
  \[
    a_s a^t
    = \mu^{k, j, n}_s m^t_{p, \ell, k}
    \sigma^p_n a^\ell a_j
  \]
  (with implied summations),
  which uniquely defines a Hopf algebra
  structure on $D(A)$.
\end{remark}

\begin{theorem}
  The Hopf algebra $D(A)$ is quasitriangular,
  with
  \[
    R = \sum_i (a_i \otimes 1_{A^*}) \otimes (1_A \otimes a^i).
  \]
\end{theorem}

\begin{example}
  Let $G$ be a finite group and $A = \C[G]$.
  Then $A^* = \mathcal{F}(G)$, and we have
  \[
    D(A) 
    = \C[G] \ltimes \mathcal{F}(G).
  \]
\end{example}
